\documentclass[11pt,a4paper]{article}
\usepackage[margin=2.5cm]{geometry}
\usepackage{amsmath,amssymb}

\title{Statistical Learning Theory\\Exercise Sheet 1}
\author{Group 9 - Herzog, Redlich, Kriewald}
\date{}

\begin{document}
\maketitle

\section*{Discussion: II(c)}

Within the training range $[1,5]$, the four networks have similar errors.
They are slightly less accurate than $h_3$ at $L=1,2$, but more accurate at
$L=4,5$. At $L=3$, the differences from $h_3$ are small. A larger network
therefore does not consistently improve performance within this range.

For larger endpoints, the network errors increase sharply. At $L=6$, all
four networks have larger errors than $h_2$ and $h_3$. At $L=8$, they are
worse than all three comparison rules. Their MSE at $L=8$ ranges from about
$3.8$ to $5.9$, compared with $0.52$ for $h_3$. The learned relationship
therefore does not extend reliably to these larger endpoints.

At $L=0.5$, performance depends strongly on the architecture. The smallest
network has an MSE of about $0.004$, whereas the deepest has about $0.18$.
Poor performance outside the training range is thus not uniform across
architectures. Unlike the networks, $h_1,h_2,h_3$ are not fitted to the
training interval. Of these three rules, $h_3$ is the most accurate at
every tested endpoint.

\end{document}
